\clearpage
\fancyhead[L]{\small COMET.}
\hypertarget{native-474-lower}{}
\begin{center}
\rule{0.55\linewidth}{0.4pt}\\[2pt]
Astronomische Nachrichten. Vol. XXII. No. 521. pp. 277. 278. 28 November 1844.\\[2pt]
\rule{0.55\linewidth}{0.4pt}
\end{center}
\begin{flushright}Göttingen, 15 October and 11 November 1844.\end{flushright}

On 6 October I made a couple of comparisons of the \emph{comet} with Ceti 200 \textsc{Bode}. Only the right ascension can be derived with certainty from them.

\begin{center}
1844 October 6.\quad $10^{\rm u}\ 40^{\rm m}\ 32^{\rm s}$ mean time\quad R.A. $18^\circ\ 50'\ 51''1$
\end{center}

For the stars with which I compared the comet on 14 October, I take the following apparent positions on the day of observation, based on the various more recent meridian observations:

\begin{center}
\setlength{\tabcolsep}{3pt}
\begin{tabular}{llll}
Star 1. & \shortstack[l]{Right ascension\\$19^\circ\ 46'\ 4''0$} & \shortstack[l]{Declination\\$0^\circ\ 56'\ 51''4$} & south \\
\hfill 2. & \hfill $19\ 52\ 18.8$ & \hfill $0\ 44\ 45.5$ & \\
\end{tabular}
\end{center}
\noindent On this basis I find the comet’s position to be

\begin{center}
\setlength{\tabcolsep}{3pt}
\begin{tabular}{lll}
1844 October 14. & $9^{\rm u}\ 27^{\rm m}\ 6^{\rm s}$ mean time & R.A. $20^\circ\ 16'\ 11''3$ \\
 & $9\ 32\ 7$ & Southern decl. $0\ 44\ 30.3$ \\
\end{tabular}
\end{center}

\Needspace{7\baselineskip}
\SourceLeaf{475}
Dr. \textsc{Goldschmidt} has calculated the following elliptic elements for the comet:

\begin{center}
Passage through perihelion 1844 Sept. 2.49152 mean Berlin time.
\end{center}
{\small
\setlength{\arraycolsep}{2.5pt}
\[
\left.\begin{array}{llr}
\text{Longitude of perihelion} & \text{.\ .\ .\ .\ .} & 342^\circ\ 29'\ 44''9 \\
\text{Ascending node} & \text{.\ .\ .\ .\ .} & 63\ 48\ 55.2
\end{array}\right\}
\begin{array}{l}
\text{reckoned from the mean equinox} \\
\text{1844 Sept. 21.5}
\end{array}
\]
}
\begin{center}
\setlength{\tabcolsep}{4pt}
\begin{tabular}{llr}
Inclination of orbit & \mbox{.\ .\ .\ .\ .} & $2\ 55\ 1.9$ \\
\multicolumn{2}{l}{Logarithm of the semimajor axis} & $0.4929151$ \\
Eccentricity & \mbox{.\ .\ .\ .\ .\ .\ .\ .} & $0.6186103$ \\
\multicolumn{3}{l}{Direct motion.}
\end{tabular}
\end{center}

The sidereal period that follows from these elements is 2400.339 days. The elements reproduce exactly the 3 observations on which they are based (Aug. 23, Sept. 21, Oct. 15). Dr. \textsc{Goldschmidt} will soon compare all the observations that have become known with these elements. For now he has compared the three observations of 14 Oct., 1 Nov. and 2 Nov., obtaining the following differences:

\begin{center}
\begin{tabular}{lrrl}
 & in R.A. & in declination & \\
Oct. 14 & $-0''9$ & $-8''3$ & Gauss \\
Nov. 1 & $+8.7$ & $-10.3$ & Petersen \\
\hfill 2 & $-12.3$ & $-8.4$ & Rümker \\
\end{tabular}
\end{center}

The following ephemeris was calculated from these elements.

